% ============================================================
% Section 92: 一维紧束缚能带、态密度与热力学性质
% ============================================================
\section{固体理论：一维紧束缚能带、态密度与热力学性质}
\label{sec:1d-tight-binding-thermodynamics}

\begin{modelbox}[题目：一维紧束缚模型的能带、态密度与热力学性质]
相距为 $b$ 的原子规则组成一维金属，每个原子有一个价电子。使用紧束缚近似，原子波函数为 $\Psi(\bm{r}-\bm{R}_i)$，$\bm{R}_i=i\,b\hat{\bm{x}}$。假定
\begin{align*}
\int\Psi^*(\bm{r}-\bm{R}_i)H\Psi(\bm{r}-\bm{R}_i)\,d^3r&=-E_0,\\
\int\Psi^*(\bm{r}-\bm{R}_i)H\Psi(\bm{r}-\bm{R}_{i\pm1})\,d^3r&=-V,\\
\int\Psi^*(\bm{r}-\bm{R}_i)H\Psi(\bm{r}-\bm{R}_{i\pm j})\,d^3r&=0\quad(j\ge2).
\end{align*}
计算：
\begin{enumerate}
    \item 电子能带结构 $E(k)$；
    \item 电子态密度；
    \item 电子对结合能的贡献（相对于 $-E_0$）；
    \item 电子对比热的贡献（$k_BT\ll V$）。
\end{enumerate}
\end{modelbox}

\begin{tipbox}[考点快速分析]
\begin{itemize}
    \item \textbf{紧束缚能带}：最近邻跃迁给出余弦色散 $E(k)=-E_0-2V\cos kb$
    \item \textbf{一维态密度}：$g(E)\propto|dE/dk|^{-1}$，带边有范霍夫奇点
    \item \textbf{半填充}：每个原子1电子，$k_F=\pi/(2b)$，$E_F=-E_0$
    \item \textbf{低温比热}：索末菲展开，$c_e\propto T$
\end{itemize}
\end{tipbox}

\begin{derivationbox}[标准解题过程]
\textbf{Step 1: 电子能带结构}

布洛赫波函数
\begin{equation}
\varphi_k(x)=\frac{1}{\sqrt{N}}\sum_n e^{iknb}\Psi(x-nb).
\end{equation}

能量本征值
\begin{align}
E(k)&=\langle\varphi_k|H|\varphi_k\rangle \\
&=-E_0-V(e^{ikb}+e^{-ikb}) \\
&=\boxed{-E_0-2V\cos(kb),\qquad -\frac{\pi}{b}<k\le\frac{\pi}{b}.}
\end{align}

\textbf{Step 2: 电子态密度}

一维晶格中，计入自旋简并，$k$ 空间态密度 $L/\pi$（$L=Nb$）。由
\begin{equation}
\frac{dE}{dk}=2Vb\sin(kb),
\end{equation}
得
\begin{equation}
g(E)=\frac{L}{\pi}\Bigl|\frac{dk}{dE}\Bigr|=\frac{N}{\pi}\frac{1}{2V|\sin(kb)|}.
\end{equation}

利用 $\cos(kb)=-(E+E_0)/(2V)$，$\sin(kb)=\sqrt{1-\cos^2(kb)}$，得
\begin{equation}
\boxed{g(E)=\frac{N}{\pi V}\frac{1}{\sqrt{1-\bigl(\frac{E+E_0}{2V}\bigr)^2}},\qquad -E_0-2V\le E\le-E_0+2V.}
\end{equation}
在带边 $E=-E_0\pm2V$ 处，$g(E)\to\infty$，出现\textbf{范霍夫奇点}。

\textbf{Step 3: 电子对结合能的贡献}

每个原子贡献1电子，总电子数 $N$。能带关于 $k=0$ 对称，半填充时费米波矢
\begin{equation}
k_F=\frac{\pi}{2b}.
\end{equation}

总电子能量
\begin{align}
E_{\text{tot}}&=\int_{-k_F}^{k_F}E(k)\frac{L}{\pi}\,dk \\
&=\frac{Nb}{\pi}\int_{-\pi/2b}^{\pi/2b}\bigl[-E_0-2V\cos(kb)\bigr]dk \\
&=N\Bigl(-E_0-\frac{4V}{\pi}\Bigr).
\end{align}

每个电子平均能量 $\bar E=-E_0-4V/\pi$。相对于孤立原子能量 $(-E_0)$，每个电子对结合能的贡献为
\begin{equation}
\boxed{\Delta E=\bar E-(-E_0)=-\frac{4V}{\pi}.}
\end{equation}
（负号表示晶体中电子能量更低，体系更稳定。）

\textbf{Step 4: 电子对比热的贡献（$k_BT\ll V$）}

低温下利用索末菲展开，内能变化
\begin{equation}
\Delta E_{\text{tot}}(T)\approx\frac{\pi^2}{6}g(E_F)(k_BT)^2.
\end{equation}

半填充时 $k_F=\pi/(2b)$，$E_F=E(k_F)=-E_0$。费米面处态密度
\begin{equation}
g(E_F)=\frac{N}{\pi V}\frac{1}{\sqrt{1-0}}=\frac{N}{\pi V}.
\end{equation}

每个电子对比热的贡献
\begin{equation}
c_e=\frac{1}{N}\frac{d(\Delta E_{\text{tot}})}{dT}
=\frac{\pi^2}{6}\frac{1}{\pi V}2k_B^2T
=\boxed{\frac{\pi k_B^2}{3V}\,T.}
\end{equation}
与温度成正比，符合金属自由电子气的普遍规律。
\end{derivationbox}

\begin{formulabox}[答案汇总]
\begin{center}
\begin{tabular}{ll}
\toprule
\textbf{物理量} & \textbf{表达式} \\
\midrule
(1) 能带结构 & $E(k)=-E_0-2V\cos(kb)$ \\
(2) 态密度 & $g(E)=\dfrac{N}{\pi V}\bigl/\sqrt{1-\bigl(\frac{E+E_0}{2V}\bigr)^2}$ \\
(3) 结合能贡献 & $\Delta E=-4V/\pi$（每个电子） \\
(4) 电子比热 & $c_e=\dfrac{\pi k_B^2}{3V}\,T$（每个电子） \\
\bottomrule
\end{tabular}
\end{center}
\end{formulabox}

\begin{insightbox}[物理图像]
\begin{itemize}
    \item 一维紧束缚能带的典型特征是带宽 $4V$（从 $-E_0-2V$ 到 $-E_0+2V$），以及带边的范霍夫奇点。这些奇点在碳纳米管、量子线等实际一维系统中可通过扫描隧道谱（STS）直接观测。
    \item 半填充时 $E_F=-E_0$ 恰好位于能带中心，这是电子--空穴对称性的体现。此时系统具有最大的态密度敏感性，少量掺杂即可显著改变费米面位置。
    \item 结合能 $|\Delta E|=4V/\pi$ 反映了电子离域化（通过跃迁积分 $V$）带来的能量降低，是金属键合的本质来源。$V$ 越大，电子巡游性越强，结合能越大。
    \item 低温比热与 $T$ 线性正比，系数反比于 $V$。这表明跃迁越强（能带越宽），费米面态密度越低，电子热容越小——与三维自由电子气的行为一致。
\end{itemize}
\end{insightbox}
